\documentclass[../../main.tex]{subfiles}
\begin{document}

{\bf [解答]}

与えられた数列
\begin{align}
  Y_1=X_1, \quad Y_k=X_k+\dfrac{1}{Y_{k-1}} \quad (X_k=1,2,\cdots,6) \label{eq:1}
\end{align}
が
\begin{align}
  \dfrac{1+\sqrt{3}}{2} \leqq Y_n \leqq 1+\sqrt{3} \label{eq:2}
\end{align}
を満たす確率$p_n$を考える．両端の値は
\begin{align}
  1 < \dfrac{1+\sqrt{3}}{2} < 2 <  1+\sqrt{3} < 3 \label{eq:3}
\end{align}
であることに注意する．まず\cref{eq:1}から帰納的に
\begin{align}
  Y_k \ge 1 \label{eq:4}
\end{align}
である．

まず$n=1$の時は\cref{eq:1,eq:3}を\cref{eq:2}に代入して
\begin{align}
  X_1 = 2
\end{align}
であれば良いから，このようになる確率$p_1$は
\begin{align}
  p_1 = \dfrac{1}{6}  \label{eq:5}
\end{align}
である．

そこで$n \ge 2$の場合を考える．\cref{eq:1}を\cref{eq:2}に代入して
\begin{align}
  & \dfrac{1+\sqrt{3}}{2} \leqq X_n+\dfrac{1}{Y_{n-1}} \leqq 1+\sqrt{3} \\
  \therefore
  & \dfrac{1+\sqrt{3}}{2} - X_n \leqq \dfrac{1}{Y_{n-1}} \leqq 1+\sqrt{3} - X_n \label{eq:6}
\end{align}
となる．\cref{eq:4}から
\begin{align}
  1+\sqrt{3} - X_n > 0 \quad \therefore X_n = 1, 2 \quad (\because\text{\cref{eq:3}})
\end{align}
が必要である．以下$X_n$の値で場合分けして$Y_{n-1}$のとりうる範囲について調べる．

\subparagraph{$1^{\circ}$ $X_n=1$の時}

\cref{eq:6}に代入して
\begin{align}
  & \dfrac{-1+\sqrt{3}}{2} \leqq \dfrac{1}{Y_{n-1}} \leqq \sqrt{3} \\
  \therefore
  & \dfrac{1}{\sqrt{3}} \leqq Y_{n-1} \leqq 1 + \sqrt{3} \label{eq:7}
\end{align}
である．

\subparagraph{$2^{\circ}$ $X_n=2$の時}

\cref{eq:6}に代入して
\begin{align}
  & \dfrac{-3+\sqrt{3}}{2} \leqq \dfrac{1}{Y_{n-1}} \leqq -1 + \sqrt{3} \\
  \therefore
  & \dfrac{1+\sqrt{3}}{2} \leqq Y_{n-1} \quad (\because Y_{n-1}>0) \label{eq:8}
\end{align}
である．
\casesend

以上\cref{eq:7,eq:8}から，\cref{eq:6}が実現されるような$X_n,Y_{n-1}$の条件は
\begin{align}
  &
  \begin{cases}
    \dfrac{1}{\sqrt{3}} \leqq Y_{n-1} \leqq 1 + \sqrt{3} \text{の時}: X_n = 1 \\
    \dfrac{1+\sqrt{3}}{2} \leqq Y_{n-1} \text{の時}: X_n = 2
  \end{cases} \\
  \therefore
  \begin{cases}
    \dfrac{\sqrt{3}}{3}\le Y_{n-1}<\dfrac{1+\sqrt{3}}{2}\ \text{の時，}X_{n}=1，\\
    \dfrac{1+\sqrt{3}}{2}\le Y_{n-1}\le1+\sqrt{3}\ \text{の時，}X_{n}=1\ \text{or}\ 2，\\
    1+\sqrt{3}<Y_{n-1}\ \text{の時，}X_{n}=2
  \end{cases}
\end{align}
である．\cref{eq:4}からこれで全ての$Y_n$のとりうる値は尽くされているから，それぞれの場合の確率は
\begin{align}
  \begin{cases}
    \dfrac{1+\sqrt{3}}{2}\le Y_{n-1}\le1+\sqrt{3} \ \text{の時}: p_{n-1} \\
    \dfrac{\sqrt{3}}{3}\le Y_{n-1}<\dfrac{1+\sqrt{3}}{2}, 1+\sqrt{3}<Y_{n-1} \ \text{の時}: 1-p_{n-1}
  \end{cases}
\end{align}
となって$p_{n-1}$と$p_n$の関係は
\begin{align}
  p_{n}
  &=\dfrac{2}{6}p_{n-1}+\dfrac{1}{6}(1-p_{n-1})\\
  &=\dfrac{1}{6}p_{n-1}+\dfrac{1}{6} \\
  \therefore
  p_{n}-\dfrac{1}{5} = \dfrac{1}{6}\left(p_{n-1}-\dfrac{1}{5}\right)
\end{align}
となる．\cref{eq:5}のもとでこの漸化式を解いて，求める確率は
\begin{align}
  p_n=\dfrac{1}{5}\left\{1-\left(\dfrac{1}{6}\right)^n\right\}
\end{align}
である．

\newpage

\end{document}
